% संपूर्ण मराठी ग्रंथातील प्रकरण 15: प्रथम-क्रम तर्कशास्त्राची विन्यासमीमांसा
% हा संपादनयोग्य स्रोतखंड आहे. संकलनासाठी संपूर्ण स्रोतसंग्रहातील
% अक्षररूपे, पूर्वभाग, चिन्हव्याख्या आणि आकृती-स्रोत आवश्यक आहेत.
% मूळ: Open Logic Project; मजकूर आणि रूपांतर CC BY 4.0.
% यंत्रानुवाद, दुरुस्ती आणि तपासणी: OpenAI Codex —
% GPT-5.6 Sol आणि GPT-6 Sol, दोन्ही Ultra effort.
% दुरुस्ती-पुनरावलोकन: GPT-6.1 Sol, Ultra effort.
\chapter{प्रथम-क्रम तर्कशास्त्राची विन्यासमीमांसा}\label{fol:syn::chap}









\section{परिचय}\label{fol:syn:itx:sec}

प्रथम-क्रम तर्कशास्त्राची उपपत्ती आणि अधिउपपत्ती विकसित करण्यासाठी,
आपण प्रथम त्याच्या व्यंजकांची विन्यासमीमांसा आणि चिन्हार्थमीमांसा
परिभाषित केली पाहिजे. प्रथम-क्रम तर्कशास्त्राची व्यंजके म्हणजे पदे आणि
सूत्रे. पदे चले, स्थिरांक आणि फलने
यांपासून तयार होतात. सूत्रे मात्र विधेये आणि पदे
यांपासून तयार होतात (यांतून सर्वांत लहान, ``आण्विक'' सूत्रे
तयार होतात); मग तार्किक संयोजके आणि संख्यापक वापरून आण्विक
सूत्रांपासून अधिक जटिल सूत्रे तयार करता येतात. रचनानियम
मांडण्याचे अनेक निरनिराळे मार्ग आहेत; आपण त्यांपैकी केवळ एक शक्य
मार्ग देतो. इतर पद्धती वेगळी चिन्हे निवडतील, संयोजकांचे वेगळे संच
आदिम म्हणून निवडतील आणि कंसांचा वेगळ्या रीतीने वापर करतील (किंवा
तथाकथित पोलिश संकेतपद्धतीप्रमाणे कंस अजिबात वापरणार नाहीत). तरी सर्व
दृष्टिकोनांत सामाईक बाब ही की रचनानियम पदांचा आणि सूत्रांचा
संच \emph{विगमनाने} परिभाषित करतात. हे योग्य रीतीने केल्यास,
रचनानियमांनुसार प्रत्येक व्यंजक मूलतः एकाच रीतीने तयार होऊ शकते.
त्यामुळे तयार होणारी व्यंजके \emph{एकमेव रीतीने वाचनीय} असतात. परिणामी,
त्यांची ही विगमनाधारित व्याख्या आपल्याला त्याच पद्धतीने---विगमनाधारित
व्याख्येने---त्या व्यंजकांना अर्थ देता येतो, याची खात्री देते.












\section{प्रथम-क्रम भाषा}\label{fol:syn:fol:sec}


प्रथम-क्रम तर्कशास्त्राची व्यंजके अशा मूलभूत शब्दसंग्रहापासून तयार
केली जातात ज्यात \emph{चले}, \emph{स्थिरांक},
\emph{विधेये} आणि कधीकधी \emph{फलने} असतात. त्यांपासून,
तार्किक संयोजके, संख्यापक आणि कंस व स्वल्पविराम यांसारखी विरामचिन्हे
यांच्या साहाय्याने \emph{पदे} आणि \emph{सूत्रे} तयार केली जातात.

\begin{explain}
अनौपचारिकपणे, विधेये ही गुणधर्मांची आणि संबंधांची नावे,
स्थिरांक ही स्वतंत्र वस्तूंची नावे आणि फलने ही
संगतींची नावे आहेत. एकरूपता~$\eq$ सोडल्यास हीच \emph{अतार्किक
चिन्हे} होत आणि ती मिळून एक भाषा बनते. कोणतीही प्रथम-क्रम
भाषा~$\Lang L$ तिच्या अतार्किक चिन्हांनी निर्धारित होते. सर्वांत
सामान्य बाबतीत $\Lang L$ मध्ये प्रत्येक प्रकारची अनंत चिन्हे असतात.
\end{explain}

सामान्य बाबतीत प्रथम-क्रम तर्कशास्त्रात आपण पुढील चिन्हे वापरतो:

\begin{enumerate}
\item तार्किक चिन्हे
\begin{enumerate}
\item तार्किक संयोजके:
  \startycommalist
  \ycomma $\lnot$ (नकरण)
  \ycomma $\land$ (संधियोग)
  \ycomma $\lor$ (विकल्पयोग)
  \ycomma $\lif$ (शर्त)
  \ycomma $\liff$ (द्विशर्त)
  \ycomma $\lforall$ (वैश्विक संख्यापक)
  \ycomma $\lexists$ (अस्तित्ववाची संख्यापक).
\item असत्यता~$\lfalse$ साठीचे विधानीय स्थिर.
\item सत्यता~$\ltrue$ साठीचे विधानीय स्थिर.
\item द्विस्थानी एकरूपता~$\eq$.
\item चलांचा एक गणनीय अनंत संच: $\Obj v_0$, $\Obj v_1$, $\Obj
  v_2$, \dots
\end{enumerate}
\item प्रथम-क्रम तर्कशास्त्राची \emph{मानक भाषा} बनवणारी अतार्किक चिन्हे
\begin{enumerate}
\item प्रत्येक $n>0$ साठी $n$-स्थानी विधेयांचा एक गणनीय अनंत संच: $\Obj
  A^n_0$, $\Obj A^n_1$, $\Obj A^n_2$, \dots
\item स्थिरांकाचा एक गणनीय अनंत संच: $\Obj c_0$, $\Obj c_1$, $\Obj
  c_2$, \dots.
\item प्रत्येक $n>0$ साठी $n$-स्थानी फलनांचा एक गणनीय अनंत संच:
  $\Obj f^n_0$, $\Obj f^n_1$, $\Obj f^n_2$, \dots
\end{enumerate}
\item विरामचिन्हे: (, ) आणि स्वल्पविराम.
\end{enumerate}

आपल्या बहुतेक व्याख्या आणि निष्कर्ष प्रथम-क्रम तर्कशास्त्राच्या संपूर्ण
मानक भाषेसाठी मांडले जातील. मात्र, उपयोगानुसार आपण भाषा केवळ काही
विधेये, स्थिरांक आणि फलने पुरती मर्यादितही करू शकतो.

\begin{ex}
अंकगणिताची भाषा~$\Lang L_A$ हिच्यात एक द्विस्थानी विधेय~$<$,
एक स्थिरांक~$\Obj 0$, एक एकस्थानी फलन~$\prime$ आणि दोन
द्विस्थानी फलने~$+$ व~$\times$ असतात.
\end{ex}

\begin{ex}
संचसिद्धान्ताची भाषा~$\Lang L_Z$ हिच्यात केवळ एक द्विस्थानी
विधेय~$\in$ असते.
\end{ex}

\begin{ex}
क्रमांची भाषा~$\Lang L_\le$ हिच्यात केवळ द्विस्थानी
विधेय~$\le$ असते.
\end{ex}

पुन्हा, हे संकेत आहेत: औपचारिकदृष्ट्या ही केवळ पर्यायी नावे आहेत;
उदा., $<$, $\in$ आणि $\le$ ही $\Obj A^2_0$ ची, $\Obj 0$ हे
$\Obj c_0$ चे, $\prime$ हे $\Obj f^1_0$ चे, $+$ हे $\Obj f^2_0$ चे आणि $\times$ हे
$\Obj f^2_1$ चे पर्यायी नाव आहे.

.





\begin{intro}
वर आपण वापरलेल्या संज्ञांपेक्षा आणि चिन्हांपेक्षा वेगळ्या संज्ञा व
चिन्हे तुम्हाला परिचित असू शकतात. तर्कशास्त्राच्या ग्रंथांत (आणि
अध्यापनात) ``नकरण'' यासाठी सामान्यतः $\sim$, $\neg$ किंवा~!\ यांपैकी
एखादे, तर ``संधियोग'' यासाठी $\wedge$, $\cdot$ किंवा $\&$ यांपैकी
एखादे चिन्ह वापरतात. ``सशर्त (संकेतार्थक)'' किंवा ``अभिव्यंजन'' यांसाठी
$\rightarrow$, $\Rightarrow$ आणि $\supset$ ही चिन्हे सामान्यतः वापरली जातात.
``द्विसशर्त,'' ``द्वि-अभिव्यंजन'' किंवा
  ``(वास्तविक) सममूल्यता'' यांसाठीची चिन्हे $\leftrightarrow$,
  $\Leftrightarrow$ आणि $\equiv$ ही आहेत.
$\lfalse$ या चिन्हाला निरनिराळ्या ठिकाणी
  ``असत्यता,'' ``फाल्सम,'' ``विसंगती'' किंवा ``तळ'' म्हणतात.
$\ltrue$ या चिन्हाला निरनिराळ्या ठिकाणी
  ``सत्यता,'' ``वेरम'' किंवा ``शिखर'' म्हणतात.

स्थिरांक (ज्यांना कधीकधी नावे म्हणतात) यांच्यासाठी लॅटिन
वर्णमालेच्या सुरुवातीची लघु अक्षरे (उदा., $a$, $b$, $c$) आणि
चले साठी तिच्या शेवटची लघु अक्षरे (उदा., $x$, $y$, $z$)
वापरण्याचा संकेत आहे. संख्यापक चले सोबत जोडले जातात, उदा.,
$x$; वैश्विक संख्यापकासाठी $\forall x$, $(\forall x)$, $(x)$, $\Pi x$,
$\bigwedge_x$ आणि अस्तित्ववाची संख्यापकासाठी $\exists x$, $(\exists
x)$, $(Ex)$, $\Sigma x$, $\bigvee_x$ ही संकेतलेखनातील रूपांतरे आहेत.
\end{intro}

\begin{explain}
आपण सर्व विधानीय कारक आणि दोन्ही संख्यापक यांना भाषेची आदिम चिन्हे
मानू शकतो. त्याऐवजी आदिम चिन्हांचा लहान संच निवडून उरलेले
संकारक परिभाषित मानू शकतो. बूलीय कारकांच्या
``सत्यता-फलनात्मकदृष्ट्या संपूर्ण'' संचांत $\{ \lnot, \lor \}$,
$\{ \lnot, \land \}$ आणि $\{ \lnot,
\lif\}$ यांचा समावेश होतो---यांपैकी प्रत्येक संच कोणत्याही एका
संख्यापकासोबत जोडल्यास अभिव्यक्तीदृष्ट्या संपूर्ण प्रथम-क्रम भाषा मिळते.

इतर दोन संकारक तुम्हाला परिचित असू शकतात: शेफर स्ट्रोक~$|$
(हे नाव हेन्री शेफर यांच्यावरून दिले आहे) आणि पिअर्स बाण~$\downarrow$,
ज्याला क्वाइनचा डॅगर असेही म्हणतात. अनुक्रमे ``नँड'' आणि ``नॉर'' हे
त्यांचे नेहमीचे वाचन दिल्यास, हे कारक स्वतःच सत्यता-फलनात्मकदृष्ट्या
संपूर्ण असतात.
\end{explain}












\section{पदे आणि सूत्र}\label{fol:syn:frm:sec}

प्रथम-क्रम भाषा~$\Lang L$ दिली की, $\Lang L$ च्या मूलभूत
शब्दसंग्रहापासून तयार होणारी व्यंजके आपण परिभाषित करू शकतो. त्यांत
विशेषतः \emph{पदे} आणि \emph{सूत्रे} यांचा समावेश होतो.

\begin{defn}[पदे]
\label{fol:syn:frm:defn:terms}
$\Lang L$ मधील \emph{पदांचा} संच~$\Trm[L]$ पुढीलप्रमाणे विगमनाने
परिभाषित केला जातो:
\begin{enumerate}
\item प्रत्येक चल हे पद आहे.
\item $\Lang L$ मधील प्रत्येक स्थिरांक हे पद आहे.
\item $f$ हे $n$-स्थानी फलन आणि $t_1$, \dots, $t_n$ ही
  पदे असतील, तर $\Atom{f}{t_1, \ldots, t_n}$ हे पद आहे.
\item
  इतर काहीही पद नाही.
\end{enumerate}
ज्या पदात एकही चले नाही ते \emph{बंद पद} होय.
\end{defn}

\begin{explain}
भाषेचे विनिर्देशन करताना स्थिरांक ही पदांपासून वेगळी चिन्हांची
श्रेणी म्हणून दिसतात; पण त्याऐवजी ती शून्यस्थानी फलने मध्ये
समाविष्ट करता आली असती. मग पदांच्या व्याख्येतील दुसरी बाब आवश्यक
राहिली नसती. $n = 0$ असेल तेव्हा $\Atom{f}{t_1, \ldots, t_n}$ याचा
अर्थ केवळ स्वतंत्र $f$ असा घ्यावा लागेल.
\end{explain}

\begin{defn}[सूत्रे]
\label{fol:syn:frm:defn:formulas}
$\Lang L$ या भाषेतील \emph{सूत्रे} चा संच~$\Frm[L]$
पुढीलप्रमाणे विगमनाने परिभाषित केला जातो:
\begin{enumerate}
\item $\lfalse$ हे आण्विक सूत्र आहे.

\item $\ltrue$ हे आण्विक सूत्र आहे.

\item $R$ हे $\Lang L$ मधील $n$-स्थानी विधेय आणि $t_1$, \dots,
  $t_n$ ही $\Lang L$ मधील पदे असतील, तर $\Atom{R}{t_1,\ldots, t_n}$ हे
  आण्विक सूत्र आहे.

\item $t_1$ आणि $t_2$ ही $\Lang L$ मधील पदे असतील, तर $\Atom{\eq}{t_1, t_2}$
  हे आण्विक सूत्र आहे.

\item $\varphi$ हे सूत्र असेल, तर $\lnot \varphi$ हे
  सूत्र आहे.

\item $\varphi$ आणि $\psi$ ही सूत्रे असतील, तर $(\varphi \land
  \psi)$ हे सूत्र आहे.

\item $\varphi$ आणि $\psi$ ही सूत्रे असतील, तर $(\varphi \lor \psi)$
  हे सूत्र आहे.

\item $\varphi$ आणि $\psi$ ही सूत्रे असतील, तर $(\varphi \lif \psi)$
  हे सूत्र आहे.

\item $\varphi$ आणि $\psi$ ही सूत्रे असतील, तर $(\varphi \liff \psi)$
  हे सूत्र आहे.

\item $\varphi$ हे सूत्र आणि $x$ हे चल असेल,
  तर $\lforall[x][\varphi]$ हे सूत्र आहे.

\item $\varphi$ हे सूत्र आणि $x$ हे चल असेल,
  तर $\lexists[x][\varphi]$ हे सूत्र आहे.

\item इतर काहीही सूत्र नाही.
\end{enumerate}
\end{defn}

\begin{explain}
पदांच्या संचाची आणि सूत्रांच्या संचाची व्याख्या या
\emph{विगमनाधारित व्याख्या} आहेत. मूलतः, आपण सूत्रांचा संच
अनंतपणे अनेक टप्प्यांत तयार करतो. आरंभीच्या टप्प्यात सर्व आण्विक सूत्रे
ही सूत्रे आहेत असे घोषित करतो; हे व्याख्येतील पहिल्या काही बाबींशी,
म्हणजे
$\ltrue$,
$\lfalse$,
$\Atom{R}{t_1,\dots,t_n}$ आणि $\Atom{\eq}{t_1,t_2}$ यांच्या बाबींशी जुळते.
म्हणून ``आण्विक सूत्र'' म्हणजे या रूपांपैकी कोणतेही सूत्र.

व्याख्येतील इतर बाबी आधीच तयार केलेल्या सूत्रांपासून नवी
सूत्रे तयार करण्याचे नियम देतात. दुसऱ्या टप्प्यात या नियमांनी
आण्विक सूत्रांपासून सूत्रे तयार करता येतात. तिसऱ्या
टप्प्यात आण्विक सूत्रे आणि दुसऱ्या टप्प्यात मिळालेली सूत्रे यांपासून
नवी सूत्रे तयार करतो, आणि ही प्रक्रिया पुढे चालू राहते. अशा एखाद्या
टप्प्यात शेवटी तयार होणारी प्रत्येक गोष्ट सूत्र असते; इतर
काहीही सूत्र नसते.
\end{explain}

संकेताप्रमाणे, आपण $\eq$ हे चिन्ह त्याच्या युक्तिवादांच्या मध्ये लिहून
कंस वगळतो: $\eq[t_1][t_2]$ हा $\Atom{\eq}{t_1,t_2}$ चा संक्षेप आहे.
तसेच, $\lnot \Atom{\eq}{t_1,t_2}$ चा संक्षेप $\eq/[t_1][t_2]$ असा
लिहितो. द्विस्थानी संयोजक~$\ast$ वापरून $\psi$, $\chi$ पासून तयार केलेले
सूत्र $(\psi \ast \chi)$ लिहिताना आपण अनेकदा सर्वांत बाहेरची कंसांची
जोडी वगळून केवळ~$\psi \ast \chi$ असे लिहू.

\begin{intro}
काही तर्कशास्त्राच्या ग्रंथांत
$\lexists[x][\varphi]$
आणि
$\lforall[x][\varphi]$
यांची गणना
सूत्रे म्हणून करण्यासाठी
चल~$x$ हे~$\varphi$ मध्ये आलेले असणे आवश्यक मानतात. ही अट न
घातल्याने कोणतीही अडचण येत नाही आणि काम अधिक सोपे होते.
\end{intro}



विशिष्ट उपयोगासाठीच्या भाषेत काम करताना आपण अनेकदा द्विस्थानी
विधेये आणि फलने त्यांच्या संबंधित पदांच्या मध्ये
लिहितो; उदा., अंकगणिताच्या भाषेत $t_1 < t_2$ आणि $(t_1 + t_2)$, तर
संचसिद्धान्ताच्या भाषेत $t_1 \in t_2$. अंकगणिताच्या भाषेतील उत्तराधिकारी
फलन तर संकेताप्रमाणे त्याच्या युक्तिवादाच्या \emph{नंतर} लिहिले
जाते:~$t'$. मात्र, औपचारिकदृष्ट्या हे अनुक्रमे
$\Atom{\Obj A^2_0}{t_1, t_2}$, $\Obj f^2_0(t_1, t_2)$,
$\Atom{\Obj A^2_0}{t_1, t_2}$ आणि $f^1_0(t)$ यांचे केवळ रूढ संक्षेप आहेत.

\begin{defn}[विन्यासात्मक अनन्यता]
$\ident$ हे चिन्ह चिन्हमालांमधील विन्यासात्मक अनन्यता व्यक्त करते;
म्हणजे $\varphi \ident \psi$ तेव्हा आणि केवळ तेव्हाच, जेव्हा $\varphi$ आणि $\psi$
या समान लांबीच्या चिन्हमाला असतात आणि प्रत्येक स्थानी त्यांच्यात तेच
चिन्ह असते.
\end{defn}

$\ident$ या चिन्हाच्या दोन्ही बाजूंना जोडणीने मिळालेल्या चिन्हमाला
येऊ शकतात. उदाहरणार्थ, $\varphi \ident (\psi \lor \chi)$ याचा अर्थ असा:
आरंभीचा कंस, चिन्हमाला~$\psi$, $\lor$ हे चिन्ह, चिन्हमाला~$\chi$ आणि
शेवटचा कंस या क्रमाने जोडून मिळणारी चिन्हमाला आणि चिन्हमाला~$\varphi$
या एकच आहेत. असे असेल, तर $\varphi$ चे पहिले चिन्ह आरंभीचा कंस आहे;
$\varphi$ मध्ये दुसऱ्या चिन्हापासून सुरू होणारी $\psi$ ही उपचिन्हमाला आहे;
तिच्यानंतर $\lor$ येते; इत्यादी बाबी आपल्याला माहीत होतात.

पदे आणि सूत्रे ही मूलभूत घटकांपासून विगमनाधारित व्याख्यांनी तयार
होत असल्यामुळे, त्यांच्याविषयीच्या गोष्टी सिद्ध करण्यासाठी पुढील विगमन
तत्त्वे वापरता येतात.

\begin{lem}[\emph{पदांवरील विगमनाचे तत्त्व}]
    \label{fol:syn:frm:lem:trmind}
    $\Lang L$ ही प्रथम-क्रम भाषा असू द्या.
    एखादा गुणधर्म~$P$ असा असेल की

    \begin{enumerate}
        \item तो प्रत्येक चल~$v$ ला लागू होतो,

        \item तो $\Lang L$ मधील प्रत्येक स्थिरांक~$a$ ला लागू होतो, आणि

        \item तो $t_1$,~\dots, $t_n$ यांना लागू होत असेल आणि $f$ हे
        $\Lang L$ मधील $n$-स्थानी फलन असेल, तर तो
        $f(t_1,\dotsc,t_n)$ लाही लागू होतो
    \end{enumerate}
    (येथे $t_1$,~\dots, $t_n$ ही $\Lang{L}$ मधील पदे आहेत असे गृहीत धरले आहे),
    तर $P$ हा $\Trm[L]$ मधील प्रत्येक पदाला लागू होतो.
\end{lem}

\begin{prob}
    \ref{fol:syn:frm:lem:trmind} सिद्ध करा.
\end{prob}

\begin{lem}[\emph{सूत्रे वरील विगमनाचे तत्त्व}]
    \label{fol:syn:frm:thm:frmind}
    $\Lang L$ ही प्रथम-क्रम भाषा असू द्या.
    एखादा गुणधर्म~$P$ सर्व आण्विक सूत्रे ला लागू होत असेल आणि तो असा असेल की

    \begin{enumerate}
        \item तो $\varphi$ ला लागू होत असेल, तर तो $\lnot \varphi$ ला
            लागू होतो;
        \item तो $\varphi$ आणि~$\psi$ ला लागू होत असेल, तर तो
            $(\varphi \land \psi)$ ला लागू होतो;
        \item तो $\varphi$ आणि~$\psi$ ला लागू होत असेल, तर तो
            $(\varphi \lor \psi)$ ला लागू होतो;
        \item तो $\varphi$ आणि~$\psi$ ला लागू होत असेल, तर तो
            $(\varphi \lif \psi)$ ला लागू होतो;
        \item तो $\varphi$ आणि~$\psi$ ला लागू होत असेल, तर तो
            $(\varphi \liff \psi)$ ला लागू होतो;
        \item तो $\varphi$ ला लागू होत असेल, तर तो
            $\lexists[x][\varphi]$ ला लागू होतो;
        \item तो $\varphi$ ला लागू होत असेल, तर तो
            $\lforall[x][\varphi]$ ला लागू होतो;
  \end{enumerate}
  (येथे $\varphi$ आणि $\psi$ ही $\Lang{L}$ मधील सूत्रे आहेत असे गृहीत धरले आहे),
  तर $P$ हा $\Frm[L]$ मधील सर्व सूत्रांना लागू होतो.
\end{lem}












\section{एकमेव वाचनीयता}\label{fol:syn:unq:sec}

\begin{explain}
आपण सूत्रांची ज्या रीतीने व्याख्या केली आहे त्यामुळे प्रत्येक
सूत्राचे \emph{एकमेव वाचन} असते; म्हणजे सूत्रे साठीच्या
आपल्या रचनानियमांनुसार ते तयार करण्याचा मूलतः एकच मार्ग आणि त्याची
``अर्थलावणी'' करण्याचा एकच मार्ग असतो. तसे नसते, तर आपल्याला संदिग्ध
सूत्रे मिळाली असती, म्हणजे एकाहून अधिक वाचने किंवा अर्थलावण्या
असलेली सूत्रे---आणि हे आपल्याला टाळायचे आहे हे स्पष्ट आहे. पण
त्याहून महत्त्वाचे म्हणजे, या गुणधर्माविना आपण पुढे देणार असलेल्या
बहुतेक व्याख्या आणि सिद्धता चालणार नाहीत.

एकमेव वाचनीयतेची खात्री न देणारे सूत्रे तयार करण्याचे सदोष नियम
दिले असते तर काय झाले असते हे पाहणे, हा मुद्दा स्पष्ट करण्याचा बहुधा
सर्वोत्तम मार्ग आहे. उदाहरणार्थ, संयोजकांच्या रचनानियमांत आपण कंस
विसरलो असतो; उदा., पुढील गोष्टीला परवानगी दिली असती:
\begin{quote}
$\varphi$ आणि $\psi$ ही सूत्रे असतील, तर $\varphi \lif \psi$ हेही तसेच आहे.
\end{quote}
$\theta$ या आण्विक सूत्रापासून सुरुवात करून यामुळे $\theta \lif \theta$ तयार करता
आले असते. यापासून, $\theta$ सोबत, $\theta \lif \theta \lif \theta$ मिळाले असते. पण
हे करण्याचे दोन मार्ग आहेत:
\begin{enumerate}
\item $\theta$ हे $\varphi$ आणि $\theta \lif \theta$ हे $\psi$ मानणे.
\item $\varphi$ हे $\theta \lif \theta$ आणि $\psi$ हे~$\theta$ मानणे.
\end{enumerate}
त्याप्रमाणे, सूत्र~$\theta \lif \theta \lif \theta$ ``वाचण्याचे'' दोन मार्ग
आहेत. ते $\psi \lif \chi$ या रूपात असून $\psi$ हे $\theta$ आणि $\chi$ हे $\theta
\lif \theta$ आहे; पण ते $\psi \lif \chi$ या रूपात \emph{देखील आहे}, जेथे $\psi$
हे $\theta \lif \theta$ आणि $\chi$ हे~$\theta$ आहे.

असे झाले, तर आपल्या व्याख्या नेहमी चालणार नाहीत. उदाहरणार्थ, एखाद्या
सूत्राचा मुख्य संकारक परिभाषित करताना आपण म्हणतो: $\psi \lif \chi$
या रूपातील सूत्रात~$\lif$ ची दर्शवलेली आस्थिति हाच मुख्य संकारक
आहे. पण $\theta \lif \theta \lif \theta$ या सूत्राची वर सांगितलेल्या दोन वेगळ्या
मार्गांनी $\psi \lif \chi$ शी जुळणी करता आली, तर एका बाबतीत $\lif$ ची
पहिली आस्थिति मुख्य संकारक मिळते आणि दुसऱ्या बाबतीत दुसरी
आस्थिति मिळते. पण मुख्य संकारक हे सूत्राचे \emph{फलन} असावे
असा आपला उद्देश आहे; म्हणजे प्रत्येक सूत्रामध्ये मुख्य संकारक ची नेमकी एक आस्थिति असली पाहिजे.
\end{explain}

\begin{lem}
सूत्र~$\varphi$ मधील डाव्या आणि उजव्या कंसांची संख्या समान असते.
\end{lem}

\begin{proof}
$\varphi$ ज्या रीतीने तयार होते तिच्यावर विगमन करून हे सिद्ध करू. यासाठी
दोन गोष्टी आवश्यक आहेत: (a) प्रथम, सर्व आण्विक सूत्रांना संबंधित
गुणधर्म लागू होतो हे सिद्ध करावे लागते (विगमनाचा आधार). (b) नंतर,
दिलेल्या सूत्रांपासून नवी सूत्रे तयार केली आणि जुन्या सूत्रांना तो
गुणधर्म लागू होत असेल, तर नव्या सूत्रांनाही तो लागू होतो हे सिद्ध
करावे लागते.

$l(\varphi)$ ही $\varphi$ मधील डाव्या कंसांची संख्या आणि $r(\varphi)$ ही उजव्या
कंसांची संख्या असू द्या; त्याचप्रमाणे $l(t)$ आणि $r(t)$ या पद~$t$
मधील अनुक्रमे डाव्या आणि उजव्या कंसांच्या संख्या असू द्या.

\begin{prob}
    कोणत्याही पद~$t$ साठी $l(t) = r(t)$ हे सिद्ध करा.
\end{prob}

\begin{enumerate}
\item \indcase{\varphi}{\lfalse}{$\indfrm$ मध्ये $0$ डावे आणि $0$
    उजवे कंस आहेत.}

\item \indcase{\varphi}{\ltrue}{$\indfrm$ मध्ये $0$ डावे आणि $0$
    उजवे कंस आहेत.}

\item \indcase{\varphi}{\Atom{R}{t_1,\dots,t_n}}{$l(\indfrm) = 1 + l(t_1) +
  \dots + l(t_n) = 1 + r(t_1) + \dots + r(t_n) = r(\indfrm)$. येथे
  कोणत्याही पद~$t$ साठी $l(t) = r(t)$ ही, सरावासाठी सोडलेली, बाब
  वापरली आहे.}

\item \indcase{\varphi}{\eq[t_1][t_2]}{$l(\indfrm) = l(t_1) + l(t_2) =
  r(t_1) + r(t_2) = r(\indfrm)$.}

\item \indcase{\varphi}{\lnot \psi}{विगमन गृहीतकानुसार,
    $l(\psi) = r(\psi)$. म्हणून $l(\indfrm) = l(\psi) = r(\psi) =
    r(\indfrm)$.}

\item \indcase{\varphi}{(\psi \ast \chi)}{विगमन गृहीतकानुसार, $l(\psi) =
  r(\psi)$ आणि $l(\chi) = r(\chi)$. म्हणून $l(\indfrm) = 1 + l(\psi) + l(\chi) =
  1 + r(\psi) + r(\chi) = r(\indfrm)$.}

\item \indcase{\varphi}{\lforall[x][\psi]}{विगमन
    गृहीतकानुसार, $l(\psi) = r(\psi)$. म्हणून $l(\indfrm) = l(\psi) = r(\psi) =
    r(\indfrm)$.}





\item \indcase{\varphi}{\lexists[x][\psi]}{त्याचप्रमाणे.}
\end{enumerate}
\end{proof}

\begin{defn}[उचित पूर्वखंड]
$\psi$ ही चिन्हमाला आणि एक रिकामी नसलेली चिन्हमाला यांची जोडणी केल्यावर
चिन्हमाला~$\varphi$ मिळत असेल, तर चिन्हमाला $\psi$ ही चिन्हमाला~$\varphi$ चा
\emph{उचित पूर्वखंड} आहे.
\end{defn}

\begin{lem}\label{fol:syn:unq:lem:no-prefix}
$\varphi$ हे सूत्र आणि $\psi$ हा $\varphi$ चा उचित पूर्वखंड असेल, तर $\psi$
हे सूत्र नाही.
\end{lem}

\begin{proof}
सराव.
\end{proof}

\begin{prob}
\ref{fol:syn:unq:lem:no-prefix} सिद्ध करा.
\end{prob}

\begin{prop}
\label{fol:syn:unq:prop:unique-atomic}
$\varphi$ हे आण्विक सूत्र असेल, तर पुढीलपैकी एक आणि फक्त एकच अट त्याला
लागू होते.
\begin{enumerate}
\item $\varphi \ident \lfalse$.
\item $\varphi \ident \ltrue$.
\item $\varphi \ident \Atom{R}{t_1,\dots,t_n}$, जेथे $R$ हे $n$-स्थानी
  विधेय, $t_1$, \dots, $t_n$ ही पदे आणि $R$, $t_1$, \dots,
  $t_n$ यांपैकी प्रत्येक एकमेव रीतीने निर्धारित आहे.
\item $\varphi \ident \eq[t_1][t_2]$, जेथे $t_1$ आणि $t_2$ ही एकमेव रीतीने
  निर्धारित पदे आहेत.
\end{enumerate}
\end{prop}

\begin{proof}
सराव.
\end{proof}

\begin{prob}
\ref{fol:syn:unq:prop:unique-atomic} सिद्ध करा (सूचना: पदांसाठी
\ref{fol:syn:unq:lem:no-prefix} ची एक आवृत्ती मांडा आणि सिद्ध करा.)
\end{prob}

\begin{prop}[एकमेव वाचनीयता]
प्रत्येक सूत्राला पुढीलपैकी एक आणि फक्त एकच अट लागू होते.
\begin{enumerate}
\item $\varphi$ हे आण्विक आहे.

\item $\varphi$ हे $\lnot \psi$ या रूपात आहे.

\item $\varphi$ हे $(\psi \land \chi)$ या रूपात आहे.

\item $\varphi$ हे $(\psi \lor \chi)$ या रूपात आहे.

\item $\varphi$ हे $(\psi \lif \chi)$ या रूपात आहे.

\item $\varphi$ हे $(\psi \liff \chi)$ या रूपात आहे.

\item $\varphi$ हे $\lforall[x][\psi]$ या रूपात आहे.

\item $\varphi$ हे $\lexists[x][\psi]$ या रूपात आहे.
\end{enumerate}
शिवाय, प्रत्येक बाबतीत $\psi$, किंवा $\psi$ आणि $\chi$, एकमेव रीतीने
निर्धारित असतात. म्हणजे, उदाहरणार्थ, $\psi$, $\chi$ आणि $\psi'$, $\chi'$ अशा
दोन वेगळ्या जोड्या अस्तित्वात नाहीत की $\varphi$ हे एकाच वेळी
$(\psi \lif \chi)$ आणि $(\psi' \lif \chi')$ या दोन्ही रूपांत असेल.
\end{prop}

\begin{proof}
रचनानियमांनुसार, सूत्र आण्विक नसेल, तर त्याची सुरुवात आरंभीच्या
कंसाने~(, $\lnot$ ने, किंवा एखाद्या संख्यापकाने झाली
पाहिजे. दुसरीकडे, पुढीलपैकी एखाद्या चिन्हाने सुरू होणारे प्रत्येक
सूत्र आण्विक असले पाहिजे: विधेय, फलन,
स्थिरांक, $\lfalse$, $\ltrue$.

म्हणून आपल्याला प्रत्यक्षात एवढेच दाखवायचे आहे: $\varphi$ हे $(\psi \ast
\chi)$ या रूपात आणि $(\psi' \mathbin{\ast'} \chi')$ या रूपातही असेल, तर
$\psi \ident \psi'$, $\chi \ident \chi'$ आणि $\ast = {\ast'}$.

म्हणून $\varphi \ident (\psi \ast \chi)$ आणि $\varphi \ident (\psi'
\mathbin{\ast'} \chi')$ हे दोन्ही गृहीत धरा. मग एकतर $\psi \ident \psi'$
किंवा तसे नाही. तसे असेल, तर स्पष्टपणे $\ast = {\ast'}$ आणि $\chi
\ident \chi'$, कारण त्या स्थितीत त्या $\varphi$ च्या एकाच स्थानी सुरू होणाऱ्या
आणि समान लांबीच्या उपचिन्हमाला आहेत. उरलेली बाब म्हणजे $\psi
\not\ident \psi'$. $\psi$ आणि $\psi'$ या दोन्ही $\varphi$ च्या एकाच स्थानी सुरू
होणाऱ्या उपचिन्हमाला असल्यामुळे, त्यांपैकी एक दुसरीचा उचित पूर्वखंड
असला पाहिजे. पण \ref{fol:syn:unq:lem:no-prefix} नुसार हे अशक्य आहे.
\end{proof}













\section{सूत्राचा मुख्य संकारक}\label{fol:syn:mai:sec}

\begin{explain}
सूत्र~$\varphi$ तयार करताना सर्वांत शेवटी वापरलेल्या कारकाविषयी
बोलणे अनेकदा उपयुक्त असते. या कारकाला $\varphi$ चा \emph{मुख्य कारक}
म्हणतात. अंतःप्रेरणेने तो $\varphi$ चा ``सर्वांत बाहेरचा'' कारक होय.
उदाहरणार्थ, $\lnot \varphi$ चा मुख्य कारक $\lnot$, $(\varphi \lor \psi)$ चा
मुख्य कारक $\lor$, इत्यादी.
\end{explain}


\begin{defn}[मुख्य संकारक]
\label{fol:syn:mai:def:main-op}
सूत्र~$\varphi$ चा \emph{मुख्य संकारक} पुढीलप्रमाणे परिभाषित केला जातो:
\begin{enumerate}
\item \indcase*{\varphi}{\varphi}{$\indfrm$ ला मुख्य संकारक नाही.}

\item \indcase{\varphi}{\lnot \psi}{$\indfrm$ चा मुख्य संकारक
  हा~$\lnot$ आहे.}

\item \indcase{\varphi}{(\psi \land \chi)}{$\indfrm$ चा
  मुख्य संकारक हा~$\land$ आहे.}

\item \indcase{\varphi}{(\psi \lor \chi)}{$\indfrm$ चा
  मुख्य संकारक हा~$\lor$ आहे.}

\item \indcase{\varphi}{(\psi \lif \chi)}{$\indfrm$ चा
  मुख्य संकारक हा~$\lif$ आहे.}

\item \indcase{\varphi}{(\psi \liff \chi)}{$\indfrm$ चा
  मुख्य संकारक हा~$\liff$ आहे.}

\item \indcase{\varphi}{\lforall[x][\psi]}{$\indfrm$ चा
  मुख्य संकारक हा~$\lforall$ आहे.}

\item \indcase{\varphi}{\lexists[x][\psi]}{$\indfrm$ चा
  मुख्य संकारक हा~$\lexists$ आहे.}
\end{enumerate}
\end{defn}

प्रत्येक बाबतीत सूत्रातील मुख्य संकारक ची नेमकी दर्शवलेली
\emph{आस्थिति} आपल्याला अभिप्रेत आहे. उदाहरणार्थ,
$((\theta \lif \varphi_{E}) \lif (\varphi_{E} \lif \theta))$ हे सूत्र $(\psi \lif \chi)$ या रूपात
असून $\psi$ हे $(\theta \lif \varphi_{E})$ आणि $\chi$ हे $(\varphi_{E} \lif \theta)$ असल्यामुळे,
$\lif$ ची दुसरी आस्थिति मुख्य संकारक आहे.

\begin{explain}
ही अशा फलनाची \emph{पुनरावर्ती} व्याख्या आहे जे सर्व अनाण्विक
सूत्रे ना त्यांच्या मुख्य संकारक च्या आस्थितीकडे पाठवते.
सूत्रांची व्याख्या विगमनाने केलेली असल्यामुळे प्रत्येक
सूत्र~$\varphi$ ला \ref{fol:syn:mai:def:main-op} मधील एखादी बाब लागू होते.
त्यामुळे प्रत्येक अनाण्विक सूत्र~$\varphi$ साठी मुख्य संकारक
अस्तित्वात आहे याची खात्री मिळते. प्रत्येक सूत्राला यांपैकी
फक्त एकच अट लागू होते आणि प्रत्येक बाबतीत $\varphi$ ज्यांपासून तयार केले
आहे ती लहान सूत्रे एकमेव रीतीने निर्धारित असतात. त्यामुळे $\varphi$
मधील मुख्य संकारक ची आस्थिति एकमेव असते आणि म्हणून आपण फलन
परिभाषित केले आहे.
\end{explain}

सूत्रांचा मुख्य संकारक कोणते चिन्ह आहे यानुसार आपण त्यांना
\ref{fol:syn:mai:tab:main-op} मधील नावे देतो.

\begin{table}[!h]
\centering
\begin{tabular}{c | c | c}
मुख्य संकारक & सूत्राचा प्रकार & उदाहरण\\
\hline
काही नाही & आण्विक (सूत्र) &
$\lfalse$,
$\ltrue$,
$\Atom{R}{t_1, \dots, t_n}$,
$\eq[t_1][t_2]$\\
$\lnot$ & नकरण & $\lnot \varphi$ \\
$\land$ & संधियोग & $(\varphi \land \psi$) \\
$\lor$ & विकल्पयोग & $(\varphi \lor \psi$) \\
$\lif$ & शर्त & $(\varphi \lif \psi$) \\
$\liff$ & द्विशर्त & $(\varphi \liff \psi)$ \\
$\lforall[][]$ & वैश्विक (सूत्र)& $\lforall[x][\varphi]$ \\
$\lexists[][]$ & अस्तित्ववाची (सूत्र)& $\lexists[x][\varphi]$
\end{tabular}
\caption{सूत्रांचे मुख्य कारक आणि नावे}
\label{fol:syn:mai:tab:main-op}
\end{table}












\section{उपसूत्र}\label{fol:syn:sbf:sec}

\begin{explain}
दिलेल्या सूत्राचे ``घटक'' असलेल्या सूत्रे विषयी बोलणे अनेकदा
उपयुक्त असते. त्यांना आपण त्याची \emph{उपसूत्रे} म्हणतो. कोणतेही
सूत्र हे स्वतःचे उपसूत्र मानले जाते; $\varphi$ स्वतः सोडून
$\varphi$ चे इतर उपसूत्र म्हणजे \emph{उचित उपसूत्र} होय.
\end{explain}

\begin{defn}[तात्काळ उपसूत्र]
$\varphi$ हे सूत्र असेल, तर $\varphi$ ची \emph{तात्काळ उपसूत्रे}
पुढीलप्रमाणे विगमनाने परिभाषित केली जातात:
\begin{enumerate}
\item आण्विक सूत्रे ना तात्काळ उपसूत्रे नसतात.

\item \indcase{\varphi}{\lnot \psi}{$\indfrm$ चे एकमेव तात्काळ
    उपसूत्र हे~$\psi$ आहे.}

\item \indcase{\varphi}{(\psi \ast \chi)}{$\indfrm$ ची तात्काळ उपसूत्रे
  $\psi$ आणि $\chi$ आहेत ($\ast$ हे कोणतेही द्विस्थानी संयोजक आहे).}

\item \indcase{\varphi}{\lforall[x][\psi]}{$\indfrm$ चे एकमेव तात्काळ
    उपसूत्र हे~$\psi$ आहे.}

\item \indcase{\varphi}{\lexists[x][\psi]}{$\indfrm$ चे एकमेव तात्काळ
    उपसूत्र हे~$\psi$ आहे.}
\end{enumerate}
\end{defn}

\begin{defn}[उचित उपसूत्र]
$\varphi$ हे सूत्र असेल, तर $\varphi$ ची \emph{उचित उपसूत्रे}
पुढीलप्रमाणे पुनरावर्ती रीतीने परिभाषित केली जातात:
\begin{enumerate}
\item आण्विक सूत्रे ना उचित उपसूत्रे नसतात.

\item \indcase{\varphi}{\lnot \psi}{$\indfrm$ ची उचित
    उपसूत्रे म्हणजे~$\psi$ आणि $\psi$ ची सर्व उचित उपसूत्रे.}

\item \indcase{\varphi}{(\psi \ast \chi)}{$\indfrm$ ची उचित उपसूत्रे
  म्हणजे $\psi$, $\chi$, $\psi$ ची सर्व उचित उपसूत्रे आणि~$\chi$ ची
  सर्व उचित उपसूत्रे.}

\item \indcase{\varphi}{\lforall[x][\psi]}{$\indfrm$ ची उचित
    उपसूत्रे म्हणजे~$\psi$ आणि $\psi$ ची सर्व उचित
    उपसूत्रे.}

\item \indcase{\varphi}{\lexists[x][\psi]}{$\indfrm$ ची उचित
    उपसूत्रे म्हणजे~$\psi$ आणि $\psi$ ची सर्व उचित
    उपसूत्रे.}
\end{enumerate}
\end{defn}

\begin{defn}[उपसूत्र]
$\varphi$ ची उपसूत्रे म्हणजे स्वतः $\varphi$ आणि त्याची सर्व उचित
उपसूत्रे.
\end{defn}

\begin{explain}
तात्काळ उपसूत्रे आणि उचित उपसूत्रे यांच्या व्याख्यांतील
सूक्ष्म फरक लक्षात घ्या. पहिल्या बाबतीत, $\varphi$ च्या प्रत्येक शक्य
रूपासाठी आपण सूत्र~$\varphi$ ची तात्काळ उपसूत्रे थेट परिभाषित केली
आहेत. ही बाबींनुसार केलेली स्पष्ट व्याख्या आहे आणि त्या बाबी
सूत्रांच्या संचाच्या विगमनाधारित व्याख्येला समांतर आहेत. दुसऱ्या
बाबतीतही सर्व सूत्रांचा संच ज्या रीतीने परिभाषित केला आहे तिला
आपण समांतर गेलो आहोत; पण प्रत्येक बाबतीत लहान सूत्रे $\psi$, $\chi$
यांची उचित उपसूत्रे या सूत्रे ना स्वतःलाही समाविष्ट करून
घेतली आहेत.
त्यामुळे व्याख्या \emph{पुनरावर्ती} होते. सर्वसाधारणपणे, विगमनाने
परिभाषित केलेल्या संचावरील फलनाची व्याख्या (आपल्या बाबतीत हा संच म्हणजे
सूत्रे) पुनरावर्ती असते, जर फलनाच्या व्याख्येतील बाबी स्वतः त्या
फलनाचा वापर करत असतील. मात्र ती सुनिर्धारित असण्यासाठी, आपण परिभाषित
करत असलेल्या युक्तिवादाच्या ``आधी'' येणाऱ्या युक्तिवादांसाठीचीच फलनाची
मूल्ये वापरतो याची खात्री केली पाहिजे---आपल्या बाबतीत $(\psi \ast \chi)$
ची ``उचित उपसूत्र'' परिभाषित करताना आपण केवळ ``आधीच्या''
सूत्रे $\psi$ आणि $\chi$ यांची उचित उपसूत्रे वापरतो.
\end{explain}

\begin{prop}
\label{fol:syn:sbf:prop:subfrm-trans}
$\psi$ हे $\varphi$ चे उपसूत्र आणि $\chi$ हे $\psi$ चे उपसूत्र आहे असे गृहीत धरा.
मग $\chi$ हे $\varphi$ चे उपसूत्र आहे. दुसऱ्या शब्दांत, उपसूत्र-संबंध
संक्रमक आहे.
\end{prop}

\begin{prob}
\ref{fol:syn:sbf:prop:subfrm-trans} सिद्ध करा.
\end{prob}

\begin{prop}
\label{fol:syn:sbf:prop:count-subfrms}
$\varphi$ हे $n$ संयोजके आणि संख्यापक असलेले सूत्र आहे असे गृहीत धरा.
मग $\varphi$ ची जास्तीत जास्त $2n+1$ उपसूत्रे आहेत.
\end{prop}

\begin{prob}
\ref{fol:syn:sbf:prop:count-subfrms} सिद्ध करा.
\end{prob}












\section{रचनाक्रमिका}\label{fol:syn:fseq:sec}

सूत्रांची विगमनाधारित व्याख्या देणे आणि त्याला पूरक तंत्र म्हणून
सूत्रांचे गुणधर्म विगमनाने सिद्ध करणे हा सुबक व कार्यक्षम
दृष्टिकोन आहे. तथापि, सूत्रांच्या रचनेचा तळापासून वर जाणारा,
टप्प्याटप्प्याचा दृष्टिकोनही उपयुक्त ठरू शकतो; येथे आपण तो
\emph{रचनाक्रमिका} या संकल्पनेच्या साहाय्याने मांडतो.

पदे आणि सूत्रे त्यांच्या विगमनाधारित व्याख्यांचा निर्देश न करता
या रीतीने कशी मांडता येतात हे दाखवण्यासाठी, आपण प्रथम एखाद्या
भाषा~$\Lang L$ मधील चिन्हांपासून बनलेल्या स्वेच्छ चिन्हमालेची संकल्पना
मांडतो.

\begin{defn}[चिन्हमाला]
\label{fol:syn:fseq:defn:string}
$\Lang L$ ही प्रथम-क्रम भाषा आहे असे समजा. \emph{$\Lang L$-चिन्हमाला}
म्हणजे $\Lang L$ च्या चिन्हांची सांत क्रमिका. संदर्भावरून भाषा~$\Lang L$
स्पष्टपणे ठरलेली असेल, तर $\Lang L$-चिन्हमालेला आपण अनेकदा फक्त
\emph{चिन्हमाला} म्हणू.
\end{defn}

\begin{ex}
कोणत्याही प्रथम-क्रम भाषा $\Lang L$ साठी सर्व $\Lang L$-सूत्रे
या $\Lang L$-चिन्हमाला असतात; पण याचे उलट खरे नसते. उदाहरणार्थ,
\[)(\Obj v_0\lif\lexists\] ही $\Lang L$-चिन्हमाला आहे, पण
$\Lang L$-सूत्र नाही.
\end{ex}

\begin{defn}[पदांसाठी रचनाक्रमिका]
\label{fol:syn:fseq:defn:fseq-trm}
$\Lang L$-चिन्हमालांची सांत क्रमिका $\tuple{t_0,\dotsc,t_n}$ ही पद~$t$
साठी \emph{रचनाक्रमिका} असते, जर $t \ident t_n$ आणि प्रत्येक $i \leq n$
साठी एकतर $t_i$ हे चल किंवा स्थिरांक असेल, अथवा $\Lang
L$ मध्ये एखादे $k$-स्थानी फलन~$f$ असेल आणि असे
$m_1,\dotsc,m_k < i$ अस्तित्वात असतील की
$t_i \ident f(t_{m_1},\dotsc,t_{m_k})$. फरक स्पष्ट करणे आवश्यक असेल,
तर पदांसाठीच्या रचनाक्रमिकांना आपण \emph{पद-रचनाक्रमिका} म्हणू.

\end{defn}

\begin{ex}
क्रमिका
\[    \tuple{\Obj c_0, \Obj v_0, \Atom{\Obj f^2_0}{\Obj c_0, \Obj v_0}, \Atom{\Obj f^1_0}{\Atom{\Obj f^2_0}{\Obj c_0, \Obj v_0}}}
\]
ही पद $\Atom{\Obj f^1_0}{\Atom{\Obj
f^2_0}{\Obj c_0, \Obj v_0}}$ साठी रचनाक्रमिका आहे; त्याचप्रमाणे
\[    \tuple{\Obj v_0, \Obj c_0, \Atom{\Obj f^2_0}{\Obj c_0, \Obj v_0}, \Atom{\Obj f^1_0}{\Atom{\Obj f^2_0}{\Obj c_0, \Obj v_0}}}.
\]
हीदेखील त्याची रचनाक्रमिका आहे.
\end{ex}

\begin{defn}[सूत्रांसाठी रचनाक्रमिका]
\label{fol:syn:fseq:defn:fseq-frm}
$\Lang L$-चिन्हमालांची सांत क्रमिका $\tuple{\varphi_0,\dotsc,\varphi_n}$ ही
$\varphi$ साठी \emph{रचनाक्रमिका} असते, जर $\varphi \ident \varphi_n$ आणि प्रत्येक
$i \leq n$ साठी एकतर $\varphi_i$ हे आण्विक सूत्र असेल, अथवा पुढीलपैकी
एखादी अट पूर्ण करणारे $j,k < i$ आणि एखादे चल~$x$ अस्तित्वात
असतील:
\begin{enumerate}
    \item $\varphi_i \ident \lnot \varphi_j$.
    \item $\varphi_i \ident (\varphi_j \land \varphi_k)$.
    \item $\varphi_i \ident (\varphi_j \lor \varphi_k)$.
    \item $\varphi_i \ident (\varphi_j \lif \varphi_k)$.
    \item $\varphi_i \ident (\varphi_j \liff \varphi_k)$.
    \item $\varphi_i \ident \lforall[x][\varphi_j]$.
    \item $\varphi_i \ident \lexists[x][\varphi_j]$.
\end{enumerate}
फरक स्पष्ट करणे आवश्यक असेल, तर सूत्रांसाठीच्या रचनाक्रमिकांना आपण
\emph{सूत्र-रचनाक्रमिका} म्हणू.
\end{defn}

\begin{ex}
\[\tuple{
    \Atom{\Obj A^1_0}{\Obj v_0},
    \Atom{\Obj A^1_1}{\Obj c_1},
    (\Atom{\Obj A^1_1}{\Obj c_1} \land \Atom{\Obj A^1_0}{\Obj v_0}),
    \lexists[\Obj v_0][(\Atom{\Obj A^1_1}{\Obj c_1} \land \Atom{\Obj A^1_0}{\Obj v_0})]
}
\]
ही $\lexists[\Obj v_0][(\Atom{\Obj A^1_1}{\Obj c_1} \land \Atom{\Obj
A^1_0}{\Obj v_0})]$ ची एक रचनाक्रमिका आहे; त्याचप्रमाणे
\begin{multline*}
\tuple{
    \Atom{\Obj A^1_0}{\Obj v_0},
    \Atom{\Obj A^1_1}{\Obj c_1},
    (\Atom{\Obj A^1_1}{\Obj c_1} \land \Atom{\Obj A^1_0}{\Obj v_0}),
    \Atom{\Obj A^1_1}{\Obj c_1},\\
    \lforall[\Obj v_1][\Atom{\Obj A^1_0}{\Obj v_0}],
    \lexists[\Obj v_0][(\Atom{\Obj A^1_1}{\Obj c_1} \land \Atom{\Obj A^1_0}{\Obj v_0})]
}.
\end{multline*}

दुसऱ्या उदाहरणावरून दिसते की रचनाक्रमिकेत ``अडगळ'' असू शकते:
अनावश्यक असलेली किंवा रचनेत काहीही हातभार न लावणारी सूत्रे.
\end{ex}

\begin{prop}\label{fol:syn:fseq:prop:formed}
$\Frm[L]$ मधील प्रत्येक सूत्र~$\varphi$ ला एक रचनाक्रमिका असते.
\end{prop}

\begin{proof}
$\varphi$ आण्विक आहे असे समजा. मग क्रमिका~$\tuple{\varphi}$ ही $\varphi$ साठी
रचनाक्रमिका आहे.

आता $\psi$ आणि~$\chi$ यांना अनुक्रमे $\tuple{\psi_0,\dotsc,\psi_n}$ आणि
$\tuple{\chi_0,\dotsc,\chi_m}$ या रचनाक्रमिका आहेत असे समजा.

\begin{enumerate}
  \item जर $\varphi \ident \lnot \psi$, तर
      $\tuple{\psi_0,\dotsc,\psi_n,\lnot \psi_n}$ ही $\varphi$ साठी
      रचनाक्रमिका आहे.
  \item जर $\varphi \ident (\psi \land \chi)$, तर
      $\tuple{\psi_0,\dotsc,\psi_n,\chi_0,\dotsc,\chi_m,(\psi_n \land \chi_m)}$
      ही $\varphi$ साठी रचनाक्रमिका आहे.
  \item जर $\varphi \ident (\psi \lor \chi)$, तर
      $\tuple{\psi_0,\dotsc,\psi_n,\chi_0,\dotsc,\chi_m,(\psi_n \lor \chi_m)}$
      ही $\varphi$ साठी रचनाक्रमिका आहे.
  \item जर $\varphi \ident (\psi \lif \chi)$, तर
      $\tuple{\psi_0,\dotsc,\psi_n,\chi_0,\dotsc,\chi_m,(\psi_n \lif \chi_m)}$
      ही $\varphi$ साठी रचनाक्रमिका आहे.
  \item जर $\varphi \ident (\psi \liff \chi)$, तर
      $\tuple{\psi_0,\dotsc,\psi_n,\chi_0,\dotsc,\chi_m,(\psi_n \liff \chi_m)}$
      ही $\varphi$ साठी रचनाक्रमिका आहे.
  \item जर $\varphi \ident \lforall[x][\psi]$, तर
      $\tuple{\psi_0,\dotsc,\psi_n,\lforall[x][\psi_n]}$
      ही $\varphi$ साठी रचनाक्रमिका आहे.
  \item जर $\varphi \ident \lexists[x][\psi]$, तर
      $\tuple{\psi_0,\dotsc,\psi_n,\lexists[x][\psi_n]}$
      ही $\varphi$ साठी रचनाक्रमिका आहे.
\end{enumerate}
सूत्रे वरील विगमनाच्या तत्त्वानुसार प्रत्येक सूत्राला एक
रचनाक्रमिका असते.
\end{proof}

याचे उलट विधानही आपण सिद्ध करू शकतो. ते महत्त्वाचे आहे, कारण त्यामुळे
सूत्रे परिभाषित करण्याच्या आपल्या दोन्ही पद्धती सममूल्य आहेत, म्हणजे
त्या समान परिणाम देतात, हे दिसते. तसेच रचनाक्रमिकांच्या लांबीवर साधे
विगमन वापरून सूत्रांविषयीची प्रमेये सिद्ध करता येतात.

\begin{lem}
\label{fol:syn:fseq:lem:fseq-init}
$\tuple{\varphi_0,\dotsc,\varphi_n}$ ही $\varphi_n$ साठी रचनाक्रमिका आहे आणि
$k \leq n$ असे समजा. मग $\tuple{\varphi_0,\dotsc,\varphi_k}$ ही $\varphi_k$ साठी
रचनाक्रमिका आहे.
\end{lem}

\begin{proof}
अभ्यासप्रश्न.
\end{proof}

\begin{prob}
\ref{fol:syn:fseq:lem:fseq-init} सिद्ध करा.
\end{prob}

\begin{thm}
\label{fol:syn:fseq:thm:fseq-frm-equiv}
$\Frm[L]$ हा अशा सर्व $\Lang L$-चिन्हमाला~$\varphi$ चा संच आहे की $\varphi$
साठी एखादी सूत्र-रचनाक्रमिका अस्तित्वात आहे.
\end{thm}

\begin{proof}
भाषा~$\Lang L$ मधील ज्या सर्व चिन्हमालांना रचनाक्रमिका आहे, त्यांचा
संच~$F$ असू द्या. $\Frm[L] \subseteq F$ हे
\ref{fol:syn:fseq:prop:formed} मध्ये आपण पाहिले आहे; म्हणून आता उलट
समावेश सिद्ध करू.

$\varphi$ ला $\tuple{\varphi_0,\dotsc,\varphi_n}$ ही रचनाक्रमिका आहे असे समजा.
$n$ वरील प्रबल विगमनाने $\varphi \in \Frm[L]$ हे सिद्ध करू. लांबी
$m < n$ असलेली रचनाक्रमिका ज्या प्रत्येक चिन्हमालेला आहे, ती
$\Frm[L]$ मध्ये आहे, हे आपले विगमन गृहीतक आहे. रचनाक्रमिकेच्या
व्याख्येनुसार एकतर $\varphi \ident \varphi_n$ आण्विक आहे किंवा पुढीलपैकी एखादी
अट पूर्ण करणारे $j,k < n$ अस्तित्वात असले पाहिजेत:
\begin{enumerate}
    \item $\varphi \ident \lnot \varphi_j$.
    \item $\varphi \ident (\varphi_j \land \varphi_k)$.
    \item $\varphi \ident (\varphi_j \lor \varphi_k)$.
    \item $\varphi \ident (\varphi_j \lif \varphi_k)$.
    \item $\varphi \ident (\varphi_j \liff \varphi_k)$.
    \item $\varphi \ident \lforall[x][\varphi_j]$.
    \item $\varphi \ident \lexists[x][\varphi_j]$.
\end{enumerate}
आता प्रत्येक शक्यतेचा स्वतंत्र विचार करू. $\varphi$ आण्विक असेल, तर
$\varphi_n \in \Frm[L]$. त्याऐवजी $\varphi \ident (\varphi_j \land \varphi_k)$ असे समजा.


\ref{fol:syn:fseq:lem:fseq-init} नुसार
$\tuple{\varphi_0,\dotsc,\varphi_j}$ आणि $\tuple{\varphi_0,\dotsc,\varphi_k}$ या अनुक्रमे
$\varphi_j$ आणि~$\varphi_k$ साठी रचनाक्रमिका आहेत. या $\varphi$ च्या
रचनाक्रमिकेच्या उचित आरंभीच्या उपक्रमिका असल्यामुळे दोन्हींची लांबी
$n$ पेक्षा कमी आहे. म्हणून विगमन गृहीतकानुसार $\varphi_j$ आणि~$\varphi_k$ ही
$\Frm[L]$ मध्ये आहेत; त्यामुळे सूत्राच्या व्याख्येनुसार
$(\varphi_j \land \varphi_k)$ हेही त्यात आहे. इतर शक्यता समांतर युक्तिवादाने
सिद्ध होतात.
\end{proof}

पदांसाठीच्या रचनाक्रमिकांना सूत्रांच्या रचनाक्रमिकांसारखेच गुणधर्म
असतात.

\begin{prop}
\label{fol:syn:fseq:prop:fseq-trm-equiv}
$\Trm[L]$ हा अशा सर्व $\Lang L$-चिन्हमाला $t$ चा संच आहे की $t$
साठी एखादी पद-रचनाक्रमिका अस्तित्वात आहे.
\end{prop}

\begin{proof}
अभ्यासप्रश्न.
\end{proof}

\begin{prob}
\ref{fol:syn:fseq:prop:fseq-trm-equiv} सिद्ध करा.
सूचना: \ref{fol:syn:fseq:thm:fseq-frm-equiv} च्या सिद्धतेत वापरलेली
रणनीती वापरा.
\end{prob}

रचनाक्रमिकेत दोन प्रकारची ``अडगळ'' येऊ शकते: पुनरावृत्त घटक आणि सूत्र
किंवा पद यांच्या रचनेशी असंबद्ध घटक. न्यूनतम रचनाक्रमिका विचारात घेऊन
आपण दोन्ही प्रकार दूर करू शकतो.

\begin{defn}[न्यूनतम रचनाक्रमिका]
\label{fol:syn:fseq:defn:minimal-fseq}
सूत्र~$\varphi$ साठीची रचनाक्रमिका $\tuple{\varphi_0, \ldots, \varphi_n}$ ही
$\varphi$ साठी \emph{न्यूनतम रचनाक्रमिका} असते, जर $\varphi$ साठीची प्रत्येक
दुसरी रचनाक्रमिका~$s$ घेतली असता~$s$ ची लांबी $n+1$ पेक्षा अधिक किंवा
तिच्याइतकी असेल.

त्याचप्रमाणे, पद~$t$ साठीची रचनाक्रमिका $\tuple{t_0, \ldots, t_n}$ ही
$t$ साठी \emph{न्यूनतम रचनाक्रमिका} असते, जर $t$ साठीची प्रत्येक
दुसरी रचनाक्रमिका~$s$ घेतली असता~$s$ ची लांबी $n+1$ पेक्षा अधिक किंवा
तिच्याइतकी असेल.
\end{defn}

एका सूत्राला किंवा पदाला एकापेक्षा अधिक न्यूनतम रचनाक्रमिका असू शकतात;
पण त्यांत नेमक्या त्याच चिन्हमाला असतील, हे लक्षात घ्या.

\begin{prop}
\label{fol:syn:fseq:prop:subformula-equivs}
पुढील विधाने सममूल्य आहेत:
\begin{enumerate}
    \item $\psi$ हे $\varphi$ चे उप-सूत्र आहे.
    \item $\psi$ हे $\varphi$ च्या प्रत्येक रचनाक्रमिकेत येते.
    \item $\psi$ हे $\varphi$ च्या एखाद्या न्यूनतम रचनाक्रमिकेत येते.
\end{enumerate}
\end{prop}

\begin{proof}
अभ्यासप्रश्न.
\end{proof}

\begin{prob}
\ref{fol:syn:fseq:prop:subformula-equivs} सिद्ध करा.
\end{prob}

\begin{history}
रेमंड स्मलियन यांनी त्यांच्या \emph{First-Order Logic}
(Smullyan, 1968) या पाठ्यपुस्तकात रचनाक्रमिका प्रथम मांडल्या.
रचनाक्रमिकांचे आणखी गुणधर्म Zuckerman (1973) यांनी प्रस्थापित केले.
\end{history}












\section{मुक्त चल आणि वाक्य}\label{fol:syn:fvs:sec}

\begin{defn}[चल च्या मुक्त आस्थिती]
\label{fol:syn:fvs:defn:free-occ}
सूत्र मधील चल च्या \emph{मुक्त} आस्थिती पुढीलप्रमाणे
विगमनाने परिभाषित केल्या जातात:
\begin{enumerate}
\item \indcase*{\varphi}{$\varphi$ आण्विक आहे}{$\indfrm$ मधील सर्व
  चल-आस्थिती मुक्त असतात.}

\item \indcase{\varphi}{\lnot \psi}{$\indfrm$ मधील मुक्त
  चल-आस्थिती म्हणजे नेमक्या $\psi$ मधील मुक्त आस्थिती.}

\item \indcase{\varphi}{(\psi \ast \chi)}{$\indfrm$ मधील मुक्त
  चल-आस्थिती म्हणजे $\psi$ मधील आस्थिती आणि त्यांच्यासोबत
  $\chi$ मधील आस्थिती.}

\item \indcase{\varphi}{\lforall[x][\psi]}{$\indfrm$ मधील मुक्त
  चल-आस्थिती म्हणजे $\psi$ मधील $x$ च्या आस्थिती सोडून इतर
  सर्व मुक्त आस्थिती.}

\item \indcase{\varphi}{\lexists[x][\psi]}{$\indfrm$ मधील मुक्त
  चल-आस्थिती म्हणजे $\psi$ मधील $x$ च्या आस्थिती सोडून इतर
  सर्व मुक्त आस्थिती.}
\end{enumerate}
\end{defn}

\begin{defn}[बद्ध चरे]
सूत्र~$\varphi$ मधील चल ची आस्थिती मुक्त नसेल, तर ती
\emph{बद्ध} असते.
\end{defn}

\begin{prob}
\ref{fol:syn:fvs:defn:free-occ} च्या धर्तीवर बद्ध चर-आस्थितींची
विगमनाधारित व्याख्या द्या.
\end{prob}

\begin{defn}[व्याप्ती]
सूत्र~$\varphi$ मध्ये $\lforall[x][\psi]$ ची एखाद्या उपसूत्राची
  आस्थिती असेल, तर $\varphi$ मधील $\psi$ च्या संबंधित आस्थितीला
  $\lforall[x]$ च्या संबंधित आस्थितीची \emph{व्याप्ती} म्हणतात.
  $\lexists[x]$ साठीही त्याचप्रमाणे.

$\varphi$ मधील संख्यापकाच्या
$\lforall[x]$ किंवा
    $\lexists[x]$ या आस्थितीची व्याप्ती $\psi$ असेल,
तर $\psi$ मधील $x$ च्या मुक्त आस्थिती
$\lforall[x][\psi]$ आणि
$\lexists[x][\psi]$ मध्ये बद्ध असतात. या आस्थिती
त्या उल्लेख केलेल्या संख्यापक-आस्थितीने \emph{बद्ध केल्या आहेत} असे
आपण म्हणतो.
\end{defn}

\begin{ex}
पुढील सूत्राचा विचार करा:
\[\lexists[\Obj v_0][\underbrace{\Atom{\Obj A^2_0}{\Obj v_0,\Obj v_1}}_{\psi}]
\]
$\psi$ हे $\lexists[\Obj v_0]$ ची व्याप्ती दर्शवते. हा संख्यापक $\psi$
मधील $\Obj v_0$ ची आस्थिती बद्ध करतो; पण $\Obj v_1$ ची आस्थिती बद्ध
करत नाही. म्हणून या बाबतीत $\Obj v_1$ हे मुक्त चर आहे.


अधिक गुंतागुंतीच्या सूत्र~$\varphi$ मध्ये हे कसे कार्य करते ते आता
पाहू शकतो:
\[\lforall[\Obj v_0][\underbrace{(\Atom{\Obj A^1_0}{\Obj v_0} \lif
    \Atom{\Obj A^2_0}{\Obj v_0, \Obj v_1})}_{\psi}] \lif \lexists[\Obj
  v_1][\underbrace{(\Atom{\Obj A^2_1}{\Obj v_0, \Obj v_1} \lor \lforall[\Obj v_0][\overbrace{\lnot \Atom{\Obj A^1_1}{\Obj v_0}}^{\theta}])}_{\chi}]
\]
$\psi$ ही पहिल्या $\lforall[\Obj v_0]$ ची व्याप्ती, $\chi$ ही
$\lexists[\Obj v_1]$ ची व्याप्ती आणि $\theta$ ही दुसऱ्या
$\lforall[\Obj v_0]$ ची व्याप्ती आहे. पहिले $\lforall[\Obj v_0]$ हे
$\psi$ मधील $\Obj v_0$ च्या आस्थिती बद्ध करते, $\lexists[\Obj v_1]$ हे
$\chi$ मधील $\Obj v_1$ ची आस्थिती बद्ध करते आणि दुसरे
$\lforall[\Obj v_0]$ हे $\theta$ मधील $\Obj v_0$ ची आस्थिती बद्ध करते.
$\Obj v_1$ ची पहिली आस्थिती आणि $\Obj v_0$ ची चौथी आस्थिती $\varphi$ मध्ये
मुक्त आहेत. $\Obj v_0$ ची शेवटची आस्थिती $\theta$ मध्ये मुक्त, पण $\chi$
आणि~$\varphi$ मध्ये बद्ध आहे.
\end{ex}

\begin{defn}[वाक्य]
सूत्र~$\varphi$ मध्ये चलांच्या मुक्त आस्थिती नसतील तेव्हा
आणि केवळ तेव्हाच ते \emph{वाक्य} असते.
\end{defn}














\section{आदेशन}\label{fol:syn:sub:sec}

\begin{defn}[पदातील आदेशन]
$s$ मधील~$x$ च्या प्रत्येक आस्थितीच्या जागी $t$ \emph{आदेशित करून}
मिळणारा परिणाम $\Subst{s}{t}{x}$ आपण पुढीलप्रमाणे पुनरावर्ती रीतीने
परिभाषित करतो:
\begin{enumerate}
\item \indcase{s}{c}{$\Subst{\indfrm}{t}{x}$ म्हणजे फक्त $s$.}

\item \indcase{s}{y}{$y$ हे चर असून $y \not\ident x$ असेल, तर
  $\Subst{\indfrm}{t}{x}$ म्हणजेही फक्त~$s$.}

\item \indcase{s}{x}{$\Subst{\indfrm}{t}{x}$ म्हणजे~$t$.}

\item\indcase{s}{\Atom{f}{t_1, \dots, t_n}}{$\Subst{\indfrmp}{t}{x}$
  म्हणजे $\Atom{f}{\Subst{t_1}{t}{x}, \dots,
  \Subst{t_n}{t}{x}}$.}
\end{enumerate}
\end{defn}

\begin{defn}
$\varphi$ मधील~$x$ ची कोणतीही मुक्त आस्थिती $t$ मधील एखादे चर बद्ध करणाऱ्या
संख्यापकाच्या व्याप्तीत येत नसेल, तर पद~$t$ हे $\varphi$ मध्ये $x$ साठी
\emph{आदेशनासाठी मुक्त} असते.
\end{defn}

\begin{ex} ~
\begin{enumerate}
\item $\Obj v_8$ हे $\lexists[\Obj
  v_3]\Atom{\Obj A^2_4}{\Obj v_3,\Obj v_1}$ मध्ये $\Obj v_1$ साठी
  मुक्त आहे.

\item $\Obj f^2_1(\Obj v_1, \Obj v_2)$ हे
  $\lforall[\Obj v_2]\Atom{\Obj A^2_4}{\Obj v_0,\Obj v_2}$ मध्ये
  $\Obj v_0$ साठी \emph{मुक्त नाही}.
\end{enumerate}
\end{ex}

\begin{defn}[सूत्र मधील आदेशन]
$\varphi$ हे सूत्र, $x$ हे चल आणि $t$ हे $\varphi$ मध्ये~$x$
साठी आदेशनासाठी मुक्त असलेले पद असेल, तर $\varphi$ मधील~$x$ च्या सर्व मुक्त
आस्थितींच्या जागी $t$ आदेशित करून मिळणारा परिणाम म्हणजे
$\Subst{\varphi}{t}{x}$.
\begin{enumerate}
\item \indcase{\varphi}{\lfalse}{$\Subst{\indfrm}{t}{x}$
    म्हणजे $\lfalse$.}

\item \indcase{\varphi}{\ltrue}{$\Subst{\indfrm}{t}{x}$
    म्हणजे $\ltrue$.}

\item \indcase{\varphi}{\Atom{P}{t_1,\dots,
    t_n}}{$\Subst{\indfrm}{t}{x}$ म्हणजे
    $\Atom{P}{\Subst{t_1}{t}{x}, \dots, \Subst{t_n}{t}{x}}$.}

\item \indcase{\varphi}{\eq[t_1][t_2]}{$\Subst{\indfrmp}{t}{x}$ म्हणजे
  $\Subst{t_1}{t}{x} = \Subst{t_2}{t}{x}$.}

\item \indcase{\varphi}{\lnot \psi}{$\Subst{\indfrmp}{t}{x}$
    म्हणजे $\lnot \Subst{\psi}{t}{x}$.}

\item \indcase{\varphi}{(\psi \land
    \chi)}{$\Subst{\indfrmp}{t}{x}$ म्हणजे $(\Subst{\psi}{t}{x} \land
    \Subst{\chi}{t}{x})$.}

\item \indcase{\varphi}{(\psi \lor
    \chi)}{$\Subst{\indfrmp}{t}{x}$ म्हणजे $(\Subst{\psi}{t}{x} \lor
    \Subst{\chi}{t}{x})$.}

\item \indcase{\varphi}{(\psi \lif
    \chi)}{$\Subst{\indfrmp}{t}{x}$ म्हणजे $(\Subst{\psi}{t}{x} \lif
    \Subst{\chi}{t}{x})$.}

\item \indcase{\varphi}{(\psi \liff
    \chi)}{$\Subst{\indfrmp}{t}{x}$ म्हणजे $(\Subst{\psi}{t}{x} \liff
    \Subst{\chi}{t}{x})$.}

\item
\indcase{\varphi}{\lforall[y][\psi]}{$y$ हे $x$ पेक्षा वेगळे चर असेल, तर
    $\Subst{\indfrmp}{t}{x}$ म्हणजे
    $\lforall[y][\Subst{\psi}{t}{x}]$; अन्यथा
    $\Subst{\indfrmp}{t}{x}$ म्हणजे फक्त $\indfrm$.}

\item
\indcase{\varphi}{\lexists[y][\psi]}{$y$ हे $x$ पेक्षा वेगळे चर असेल, तर
    $\Subst{\indfrmp}{t}{x}$ म्हणजे
    $\lexists[y][\Subst{\psi}{t}{x}]$; अन्यथा
    $\Subst{\indfrmp}{t}{x}$ म्हणजे फक्त $\indfrm$.}
\end{enumerate}
\end{defn}

\begin{explain}
आदेशन निष्फळ असू शकते, हे लक्षात घ्या: $\varphi$ मध्ये $x$ मुळीच येत नसेल,
तर $\Subst{\varphi}{t}{x}$ म्हणजे फक्त~$\varphi$.

$\varphi$ मध्ये~$x$ साठी $t$ हे आदेशनासाठी मुक्त असले पाहिजे हा निर्बंध पुढील
प्रकारची प्रकरणे वगळण्यासाठी आवश्यक आहे. $\varphi \ident
\lexists[y][x < y]$ आणि $t \ident y$ असेल, तर
$\Subst{\varphi}{t}{x}$ हे $\lexists[y][y < y]$ होईल. या प्रकरणात आदेशन
करताना मुक्त चर~$y$ हे संख्यापक $\lexists[y]$ ने ``ग्रहित'' केले जाते;
हा परिणाम अनिष्ट आहे. उदाहरणार्थ, $\lforall[x][\psi]$ लागू असेल तेव्हा
$\Subst{\psi}{t}{x}$ हेही लागू असावे, अशी आपली अपेक्षा असते. पण
$\lforall[x][\lexists[y][x < y]]$ चा विचार करा (येथे $\psi$ म्हणजे
$\lexists[y][x < y]$). हे, उदाहरणार्थ, नैसर्गिक संख्यांविषयी सत्य
असलेले वाक्य आहे: प्रत्येक संख्या~$x$ पेक्षा मोठी एखादी संख्या~$y$
असते. $x$ च्या जागी $y$ चे आदेशन करण्याची परवानगी दिली, तर आपल्याला
$\Subst{\psi}{y}{x} \ident \lexists[y][y < y]$ मिळेल; हे असत्य आहे.
$t$ मधील कोणतेही मुक्त चर $\varphi$ मधील संख्यापकाने अखेरीस बद्ध होणार नाही,
अशी अट घालून आपण हे टाळतो.
\end{explain}

अवजड लेखन टाळण्यासाठी आपण पुढील संकेत अनेकदा वापरतो: सूत्र~$\varphi$
मध्ये चल~$x$ मुक्तपणे येऊ शकत असेल, तर हे दाखवण्यासाठी आपण
$\varphi(x)$ असेही लिहितो. कोणते $\varphi$ आणि~$x$ अभिप्रेत आहेत हे स्पष्ट असेल
आणि $t$ हे पद असेल ($\varphi(x)$ मध्ये $x$ साठी ते मुक्त आहे असे गृहीत
धरलेले असेल), तर $\Subst{\varphi}{t}{x}$ याचे संक्षिप्त रूप म्हणून आपण
$\varphi(t)$ लिहितो. म्हणून, उदाहरणार्थ, ``$\varphi(t)$ याला
$\lforall[x][\varphi(x)]$ चा दृष्टांत म्हणतो,'' असे आपण म्हणू शकतो. याचा
अर्थ असा: $\varphi$ हे कोणतेही सूत्र, $x$ हे चल आणि $t$ हे
$\varphi$ मध्ये~$x$ साठी मुक्त असलेले पद असेल, तर
$\Subst{\varphi}{t}{x}$ हा $\lforall[x][\varphi]$ चा दृष्टांत असतो.



